\chapter{System Design}
\label{ch:design}

\section{Design Methodology}
Design follows a top-down refinement from context diagram to deployable components. UML 2.x viewpoints (structural, behavioral, deployment) are realized using TikZ for publication-quality vector graphics.

\section{High-Level Architecture}
% High-level system architecture (TikZ)
\begin{figure}[htbp]
\centering
\begin{tikzpicture}[
  node distance=1.1cm and 1.6cm,
  comp/.style={rectangle, draw=black!70, fill=blue!8, rounded corners=3pt,
    minimum width=2.8cm, minimum height=1cm, align=center, font=\small},
  cloud/.style={ellipse, draw=green!50!black, fill=green!10, minimum width=2.6cm,
    minimum height=1cm, align=center, font=\small},
  user/.style={rectangle, draw=orange!70, fill=orange!12, rounded corners=3pt,
    minimum width=2.4cm, minimum height=0.9cm, align=center, font=\small},
  arrow/.style={-{Stealth[length=2.2mm]}, thick}
]
\node[user] (farmer) {Farmer\\(Mobile Client)};
\node[comp, right=of farmer] (flutter) {Flutter\\Presentation Layer};
\node[comp, right=of flutter] (riverpod) {Riverpod\\State Management};
\node[comp, below=of riverpod] (repos) {Repository\\Layer};
\node[cloud, below=of repos] (firebase) {Firebase\\Auth, Firestore, RTDB};
\node[comp, right=of repos, yshift=0.6cm] (flask) {Flask\\Rule Engine};
\node[comp, right=of repos, yshift=-0.6cm] (ml) {FastAPI\\Random Forest};
\node[cloud, right=of flask, yshift=-0.6cm] (gemini) {Gemini\\LLM API};

\draw[arrow] (farmer) -- (flutter);
\draw[arrow] (flutter) -- (riverpod);
\draw[arrow] (riverpod) -- (repos);
\draw[arrow] (repos) -- (firebase);
\draw[arrow] (repos) -- (flask);
\draw[arrow] (repos) -- (ml);
\draw[arrow] (repos) -- (gemini);
\draw[arrow, dashed] (firebase) -- node[above, font=\scriptsize]{sensors} (repos);
\end{tikzpicture}
\caption{High-level layered architecture of the \HydroSmart{} ecosystem illustrating separation between mobile presentation, application state, data access, cloud persistence, and external intelligence services.}
\label{fig:high_level_architecture}
\end{figure}

\Cref{fig:high_level_architecture} illustrates logical layering: presentation (Flutter widgets), application state (Riverpod), repositories, and external services. Such separation satisfies NFR-05 maintainability.

\section{Deployment Architecture}
\begin{figure}[htbp]
\centering
\begin{tikzpicture}[
  box/.style={rectangle, draw, rounded corners, minimum width=3.2cm, minimum height=1.1cm, align=center, font=\small},
  arrow/.style={-{Stealth}, thick}
]
\node[box, fill=cyan!12] (device) {Android / iOS\\Flutter APK};
\node[box, fill=yellow!15, right=1.8cm of device] (render) {Render Cloud\\FastAPI + Gunicorn};
\node[box, fill=green!12, below=1.2cm of device] (fb) {Google Firebase\\Auth, Firestore, RTDB};
\node[box, fill=orange!12, right=1.8cm of render] (flaskhost) {Flask Host\\(optional)};
\node[box, fill=violet!10, above=1.2cm of render] (gemini) {Google AI\\Gemini 2.0};

\draw[arrow] (device) -- node[above, font=\scriptsize]{HTTPS} (render);
\draw[arrow] (device) -- node[left, font=\scriptsize]{SDK} (fb);
\draw[arrow] (device) -- node[above, font=\scriptsize]{REST} (flaskhost);
\draw[arrow] (device) -- node[above, font=\scriptsize]{API} (gemini);
\end{tikzpicture}
\caption{Deployment architecture showing distribution of the mobile client, managed Firebase services, containerized machine-learning inference, and optional rule-based backend.}
\label{fig:deployment_architecture}
\end{figure}


\section{Component Diagram}
\begin{figure}[htbp]
\centering
\begin{tikzpicture}[comp/.style={rectangle, draw, fill=blue!8, rounded corners, minimum width=2.5cm, minimum height=0.9cm, align=center, font=\scriptsize}]
\node[comp] (ui) {UI Components};
\node[comp, below=0.8cm of ui] (prov) {Providers};
\node[comp, below=0.8cm of prov] (repo) {Repositories};
\node[comp, right=2.8cm of repo] (fb) {Firebase SDK};
\node[comp, left=2.8cm of repo] (http) {HTTP Clients};
\draw[-{Stealth}, thick] (ui)--(prov)--(repo);
\draw[-{Stealth}, thick] (repo)--(fb);
\draw[-{Stealth}, thick] (repo)--(http);
\end{tikzpicture}
\caption{Component diagram of the Flutter client illustrating dependency direction from UI to infrastructure adapters.}
\label{fig:component_diagram}
\end{figure}

\section{Package Diagram}
The repository organizes Dart packages into \texttt{core}, \texttt{data}, \texttt{domain}, and \texttt{features/*}, enforcing acyclic dependencies from features toward domain abstractions.

\section{Behavioral Models}
\begin{figure}[htbp]
\centering
\begin{tikzpicture}[font=\small,
  msg/.style={-{Stealth}, thick},
  lifeline/.style={dashed, gray}]
\node (u) at (0,0) {User};
\node (ui) at (2.5,0) {UI};
\node (rp) at (5,0) {Repository};
\node (api) at (7.8,0) {Flask API};
\node (eng) at (10.5,0) {Engine};

\foreach \x in {0,2.5,5,7.8,10.5} {
  \draw[lifeline] (\x,0) -- (\x,-5.5);
}
\draw[msg] (0,-0.8) -- node[above]{input} (2.5,-0.8);
\draw[msg] (2.5,-1.6) -- node[above]{request} (5,-1.6);
\draw[msg] (5,-2.4) -- node[above]{POST} (7.8,-2.4);
\draw[msg] (7.8,-3.2) -- node[above]{score} (10.5,-3.2);
\draw[msg] (10.5,-4.0) -- (7.8,-4.0);
\draw[msg] (7.8,-4.6) -- (5,-4.6);
\draw[msg] (5,-5.2) -- (2.5,-5.2);
\end{tikzpicture}
\caption{Sequence diagram for single-crop recommendation via the Flask rule engine.}
\label{fig:sequence_recommendation}
\end{figure}


\begin{figure}[htbp]
\centering
\begin{tikzpicture}[font=\scriptsize, activity/.style={rectangle, draw, rounded corners, minimum width=2.6cm, minimum height=0.65cm, align=center}]
\node[activity, fill=green!10] (a) {Monitor Sensors};
\node[activity, below=0.5cm of a] (b) {Evaluate Thresholds};
\node[activity, below=0.5cm of b] (c) {Update Dashboard};
\node[activity, below=0.5cm of c] (d) {Notify User?};
\draw[-{Stealth}, thick] (a)--(b)--(c)--(d);
\end{tikzpicture}
\caption{Activity diagram for continuous sensor monitoring and dashboard refresh.}
\label{fig:activity_sensor}
\end{figure}

\begin{figure}[htbp]
\centering
\begin{tikzpicture}[font=\scriptsize,
  state/.style={rectangle, draw, rounded corners, minimum width=2.2cm, minimum height=0.7cm, align=center}]
\node[circle, fill=black, inner sep=1.5pt] (i) at (0,0) {};
\node[state, right=1.2cm of i] (unauth) {Unauthenticated};
\node[state, right=1.2cm of unauth] (auth) {Authenticated};
\node[state, below=1cm of auth] (load) {Loading Profile};
\draw[-{Stealth}, thick] (i)--(unauth);
\draw[-{Stealth}, thick] (unauth) -- node[above]{login} (auth);
\draw[-{Stealth}, thick] (auth)--(load);
\end{tikzpicture}
\caption{State diagram for Firebase authentication session lifecycle.}
\label{fig:state_auth}
\end{figure}

\section{Communication Diagram}
Repository objects communicate with singleton HTTP clients and Firebase SDK instances through narrow interfaces, minimizing coupling.

\section{Security Architecture}
Secrets (Gemini API keys) should migrate to Cloud Functions; current design stores configuration in Firestore \texttt{config/gemini\_config} with read access limited to authenticated users. RTDB rules must scope writes to device owners.

\section{Chapter Summary}
\Cref{ch:implementation} details module realization corresponding to design artifacts herein.
